Linguists use \term{transparency} for cases in which a relation can access a property through material that could have hidden it. This broad pattern is \term{mediated accessibility}; the paper defines its narrower grammatical relation, \term{access transparency}, as a comparative concept. An application passes four gates: an independently diagnosable bearer, a distinct mediator, evidence that the relation tracks the property, and an independently established concealing counterpart in the same constructional family. Verdicts are positive, non-instance, or indeterminate, and feature summation, whole-configuration construal, and external bearers come out as non-instances. English \term{number-transparent quantificational nouns} (QNs), whose agreement can track the \mention{of}-complement rather than the first noun, supply the main empirical case, with a corpus pilot giving partial prospective support and a feasibility limit; transparent free relatives supply a second illustration, positive under Kim's constructional analysis and indeterminate analysis-neutrally. Spanish and Tsez extend the protocol cross-linguistically. \term{Projectibility} is the warrant for extending observed regularities to specified new cases, but it's secondary here: the comparative concept predicts nothing by itself, and a language-particular covariance pattern supports that extension only when an observable base warrants non-trivial inferences. An identifiable \term{stabilizer} is a mechanism that would separately explain why the covariance holds. The result distinguishes a relational analysis that earns its keep from a label that restates the facts it describes.
